% !TeX program = xelatex
\documentclass{UAmathtalk}
\usepackage{setspace}
\renewcommand{\when}{12:00 p.m.}
\renewcommand{\where}{Hudson 0110}
\renewcommand{\where}{Hudson 0110.\\ \href{https://albany.zoom.us/j/81219478289?pwd=N5fy7vxKBXqlQ4mAq4KaeaAM5iAHi4.1}{Zoom ID: 812\,1947\,8289 - Passcode: 233444}}
\author{Micha\l{} Lipi\'nski}
\address{ISTA}
\urladdr{https://sites.google.com/view/michal-lipinski/}
\title{Space of Acyclic Partitions}
\date{Friday, November 06, 2026}


\begin{document}

\maketitle
%\renewcommand*{\abstractsize}{\normalsize}
\renewcommand*{\abstractsize}{\fontsize{13.5}{16.2}\selectfont}
\begin{abstract}
    {
        In this talk, I will introduce the notion of an acyclic partition, which can be interpreted either as the collection of level sets of a filtration of a Lefschetz complex (in particular, of a simplicial complex) or as a Morse decomposition of a combinatorial dynamical system. Our goal is to understand the space of all such acyclic partitions, within which one can study parameterized combinatorial dynamical systems. It turns out that the space of acyclic partitions admits an explicit polytopal structure. Additionally, we show that the space stratifies with respect to the (extended) Morse vector, which means that the number of critical points within each stratum remains constant. However, even when two filtration lie in the same stratum, so their critical points are in bijection, the associated Morse complexes need not be isomorphic in a strong sense. This observation leads to the construction of a within-stratum groupoid, which, as we conjecture, may indicate presence of a bifurcation. This is ongoing joint work with Thomas Wanner (George Mason University).    }
\end{abstract}
\end{document}



